% Source: https://artofproblemsolving.com/community/c3h3131356_2023_mbmt_team_round__montgomery_blair_math_tournament
% Stable identity: aops-topic-3131356-q114
Choose the parallel lines $ y=7 $ and $ y=-7 $, with $ O=(0,0)$, and reflect if needed so the third tangent lies to the right. Write $ A=(p,7),B=(q,-7)$ with $ p,q>0 $. Equal tangent lengths give the two segments along $ AB $ from the contact point as $ p,q $, so $ p+q=18 $. The line through $ A,B $ is tangent exactly when its distance to $ O $ is 7; equivalently $ pq=49 $. Thus $ OA^2=p^2+49=p(p+q)$, and similarly $ OB^2=q(p+q)$. Consequently $(OA+OB)^2=(p+q)(p+q+2\sqrt{pq})=18(18+14)=576 $. Hence $\boxed{OA+OB=24}$.
\begin{asy}
usepackage("amsmath");
size(7cm);
pair O=(0,0),A=(9+4sqrt(2),7),B=(9-4sqrt(2),-7);draw(circle(O,7));draw((-8,7)--(16,7)^^(-8,-7)--(16,-7));draw(A--B,blue);draw(A--O--B,dashed);label("$O$",O,W);label("$A$",A,NE);label("$B$",B,SE);label("$18$",(A+B)/2,E);
\end{asy}
